\chapter{Batas Serangan: Gradien Dikaburkan dan Evasi NIDS}

Sebuah kelompok mahasiswa menambahkan satu baris pembulatan pada masukan model, lalu
melaporkan bahwa serangan mereka sendiri gagal total: tingkat keberhasilan turun dari
tujuh belas persen menjadi nol koma nol dua persen. Angka itu masuk ke ringkasan
eksekutif, lalu ke slide, lalu ke abstrak. Dua minggu kemudian seorang penguji meminta
satu hal: buktikan bahwa gradien yang dipakai penyerang masih menunjuk arah yang benar.
Ternyata tidak. Pembulatan datar di hampir semua tempat, jadi turunannya nol, jadi
penyerang melangkah nol, jadi serangan yang gagal itu sebenarnya serangan yang tidak
pernah terjadi. Ketika penguji mengganti turunannya dengan pendekatan yang masuk akal,
tingkat keberhasilan kembali ke tujuh belas persen dalam hitungan menit.

Bab ini tentang kesalahan itu. Kesalahan pertama sederhana: pembulatan menghapus gradien,
dan menghapus gradien berarti membuat penyerang buta, bukan membuat model kuat. Kesalahan
kedua lebih menarik: ada kalanya pengaburan gradien memang menaikkan biaya serangan,
sehingga sebagian ketahanan itu nyata. Membedakan keduanya menuntut percobaan, bukan
retorika.

Paruh kedua bab membawa persoalan yang sama ke tempat yang jauh lebih berisik: sistem
pengesanan pencerobohan jaringan, atau NIDS. Di sana paket harus tetap sampai, protokol
harus tetap sah, keputusan diambil dalam milidetik, dan modelnya dilatih ulang setiap
hari. Semua itu membuat serangan evasi yang cantik di atas kertas menjadi tidak praktis
di lapangan. Tetapi ``tidak praktis'' bukan ``mustahil''.

\section{Gradien yang menghilang, dan rasa aman yang lahir darinya}
\tingkat{dasar}

Bab 29 mengajarkan dua serangan yang bertumpu pada satu benda: gradien kerugian
terhadap masukan. Serangan satu langkah mengambil arah tanda gradien dan melangkah
sejauh $\varepsilon$,

\begin{equation}
x' = x + \varepsilon\,\operatorname{sign}\bigl(\nabla_x L(x,y)\bigr),
\end{equation}

sedangkan serangan banyak langkah mengulang langkah kecil dan memproyeksikan hasilnya
kembali ke bola $\ell_\infty$ berjari-jari $\varepsilon$,

\begin{equation}
x_{t+1} = \Pi_{B(x,\varepsilon)}\Bigl(x_t - \alpha\,\operatorname{sign}
\bigl(\nabla_x L(x_t,y)\bigr)\Bigr).
\end{equation}

Keduanya bekerja karena satu alasan yang sudah muncul di bab tentang linearisasi:
untuk gangguan kecil, kerugian bisa didekati oleh fungsi linear, sehingga arah yang
paling cepat menaikkan kerugian adalah arah gradien. Kalau gradien itu benar, serangan
bekerja dengan biaya satu kali mundur. Kalau gradien itu salah, penyerang menjadi buta.
Dua kalimat itu tidak setara dengan ``kalau gradien salah, model aman''. Penyerang yang
buta masih bisa menebak, dan tebakan tanpa gradien, misalnya derau acak pada paket, tetap
bisa berhasil pada sebagian kasus. Yang hilang hanyalah jalan termurah, bukan seluruh
jalan.

\index{gradien dikaburkan}
Istilah \emph{gradien dikaburkan} menunjuk keadaan ketika
gradien yang dilihat penyerang tidak memberi informasi yang cukup tentang arah yang
sebenarnya menaikkan kerugian, bukan karena modelnya kuat, melainkan karena cara model
itu dibungkus. Athalye, Carlini, Wagner dan rekan-rekan memberi nama yang lebih tegas
untuk gejala ini: rasa aman palsu. Dalam tinjauan mereka atas sembilan pertahanan dari
sebuah konferensi pembelajaran mesin, tujuh bergantung pada gradien yang dikaburkan;
enam ditembus sepenuhnya dan satu sebagian, dengan biaya serangan yang kira-kira sama
dengan serangan biasa.

\begin{untukanak}
Bayangkan seorang guru yang tidak ingin muridnya menyontek, jadi ia menutup buku
catatan dengan lakban. Murid yang jujur memang jadi kesulitan. Tetapi murid yang
kreatif tidak berhenti: ia melihat bahwa jawaban soal nomor satu selalu lebih besar
daripada nomor dua, lalu menebak seluruh lembar dari pola itu. Lakban itu tidak
menghapus jawaban; ia hanya menghalangi satu jalan menuju jawaban.

Model yang membulatkan masukannya bekerja seperti lakban itu. Angka yang masuk ke
model dipotong ke kelipatan tertentu, misalnya dua angka di belakang koma, sehingga
perubahan kecil tidak terlihat sama sekali. Program yang mencoba mencari arah
``naikkan sedikit demi sedikit'' akan menemukan bahwa tidak ada yang berubah, lalu
menyerah dan melaporkan bahwa modelnya kebal. Laporan itu salah, dan bukan karena
programnya bodoh. Program itu hanya bertanya dengan cara yang tidak cocok.

Sekarang bagian yang paling penting untuk diingat. Model yang sama masih bisa salah
kalau angkanya diubah cukup banyak sehingga melewati batas potongan. Jadi pertanyaan
yang benar bukan ``apakah serangan kecil berhasil'', melainkan ``berapa besar perubahan
yang dibutuhkan''. Bab ini mengajarkan cara mengajukan pertanyaan itu, dan cara mengenali
laporan yang tidak pernah mengajukannya.
\end{untukanak}

Mengapa pembulatan menghapus gradien? Karena pembulatan adalah fungsi tangga. Ambil
$q>0$ dan definisikan $Q_q(t)=q\,\mathrm{round}(t/q)$, kelipatan $q$ terdekat. Pada
selang $kq-q/2 \le t < kq+q/2$ fungsi ini bernilai tetap $kq$, sehingga turunannya nol
di dalam selang dan tidak terdefinisi di titik batas. Himpunan titik batas itu berukuran
nol, jadi untuk fungsi mulus apa pun $f$, komposisi $f\circ Q_q$ punya gradien nol di
hampir semua titik.

\begin{proposisi}
Misalkan $Q_q\colon\mathbb{R}^d\to\mathbb{R}^d$ bekerja koordinat demi koordinat
seperti di atas, dan misalkan $f\colon\mathbb{R}^d\to\mathbb{R}$ terdiferensiasi
kontinu. Maka $\nabla (f\circ Q_q)(x)=0$ untuk hampir setiap $x\in\mathbb{R}^d$.
\end{proposisi}

\begin{proof}
Partisi $\mathbb{R}^d$ menjadi kotak-kotak setengah terbuka bersisi $q$, yang dibatasi
hiperbidang $x_j=k_jq\pm q/2$. Pada interior setiap kotak, $Q_q$
bernilai konstan, sehingga $f\circ Q_q$ juga konstan dan gradiennya nol. Komplemen dari
gabungan semua interior itu adalah gabungan dari hiperbidang batas, yang berukuran
nol. Jadi gradiennya nol di hampir semua titik.
\end{proof}

\begin{akibat}
Serangan satu langkah yang memakai gradien eksak dari $f\circ Q_q$ menghasilkan
$x'=x$ untuk setiap $\varepsilon>0$. Tingkat keberhasilan yang terukur karena itu
tepat nol, dan ``akurasi di bawah serangan'' yang dilaporkan sama persis dengan akurasi
bersih.
\end{akibat}

Akibat ini adalah inti seluruh bab. Pernyataan itu bukan tentang kekuatan model; ia
tentang aritmetika penyerang, dan angka nol pada tabel hasil tidak mengukur ketahanan,
melainkan ketiadaan langkah. Kalau kolom ``akurasi di bawah serangan'' identik dengan
akurasi bersih sampai dua angka di belakang koma, Anda bukan melihat pertahanan yang
berhasil, melainkan penyerang yang berhenti.

\begin{bacakeras}
Dua suara. Yang pertama membawa peta, yang kedua membawa senter.

--- Lihat. Petanya sudah kubuat kabur. Semua jalan menanjak kuhapus, tidak ada panah,
tidak ada angka.
--- Kalau begitu gunungnya hilang.
--- Bukan. Gunungnya masih di sana.
--- Jadi yang hilang cuma gambarnya?
--- Yang hilang cuma gambarnya.
--- Orang yang mencium bau tanah masih bisa naik.
--- Masih bisa.
--- Orang yang melempar kerikil lalu mendengar pantulannya juga masih bisa.
--- Juga masih bisa.
--- Berarti petamu bukan tembok.
--- Petamu bukan tembok. Petamu hanya peta yang menolak menjawab.

Ulangi baris terakhir tiga kali, makin pelan: peta yang menolak menjawab bukan tembok.
Bukan tembok. Bukan tembok.
\end{bacakeras}

Athalye dan rekan-rekan membagi mekanismenya menjadi tiga keluarga. Keluarga pertama,
gradien hancur, muncul ketika bagian dari model tidak terdiferensiasi atau
terdiferensiasi tetapi tidak berguna, seperti pembulatan, penyaringan median, dan
kompresi. Keluarga kedua, gradien acak, muncul ketika model mengacak sesuatu pada setiap
panggilan, sehingga gradien satu panggilan tidak berlaku untuk panggilan berikutnya.
Keluarga ketiga, gradien hilang atau meledak, muncul ketika pertahanan menekan gradien
saat pelatihan, sehingga gradien menjadi terlalu kecil atau terlalu besar untuk dipakai
melangkah.

Ketiganya punya satu ciri bersama: ia membuat pengukuran terlihat lebih baik tanpa
membuat model lebih baik, dan karena pengukuran biasanya satu-satunya bukti yang
dilaporkan, kesalahan ini menyebar cepat. Bab ini tidak mengatakan bahwa semua pertahanan
seperti itu tidak berguna; sebagian menaikkan biaya serangan, dan biaya itu nyata. Yang
salah adalah menyimpulkan ketahanan dari pengukuran yang rusak.

\section{Tujuh cara mengaburkan gradien}
\tingkat{menengah}

Tabel berikut merangkum tujuh teknik yang paling sering muncul, keluarga mekanismenya,
dan gejala yang seharusnya membuat Anda curiga. Angka pada kolom terakhir bukan
diagnosis, melainkan pemicu pertanyaan.

\begin{center}
\small
\begin{tabular}{@{}p{0.6cm}p{4.6cm}p{2.1cm}p{6.4cm}@{}}
\hline
No & Teknik & Keluarga & Gejala yang terlihat \\
\hline
1 & Kuantisasi masukan atau keluaran & hancur & gradien nol; akurasi serangan sama dengan akurasi bersih \\
2 & Praproses tak terdiferensiasi: median, blur, kompresi & hancur & gradien nol, atau melonjak tanpa pola \\
3 & Penyaringan dan penolakan sampel & hancur & keputusan diskrit; tidak ada gradien \\
4 & Transformasi acak: ubah ukuran, bingkai, geser & acak & gradien berbeda tiap panggilan \\
5 & Pengklasifikasi acak dan ansambel acak & acak & satu sampel menunjuk arah salah \\
6 & Gradien ditekan saat pelatihan & hilang & gradien kecil padahal kerugian besar \\
7 & Pertahanan berbasis optimisasi banyak langkah & hilang & gradien meledak; serangan tidak konvergen \\
\hline
\end{tabular}
\end{center}

\index{kuantisasi}
Teknik pertama, kuantisasi, sudah dibuktikan di atas, dan efeknya eksak: bukan
kira-kira nol, tetapi nol. Teknik kedua bekerja dengan cara yang sama tetapi lebih
kotor. Penyaringan median mengganti setiap nilai dengan median lingkungannya, sehingga
fungsinya tetap datar di sebagian besar tempat, tetapi batas antar potongannya tidak lagi
berjarak seragam. Kompresi seperti JPEG membuang komponen frekuensi tinggi; turunannya
ada tetapi tidak stabil, dan gradien beda hingga bisa berubah tanda ketika langkahnya
digeser. Gejalanya khas: gradien bukan nol, melainkan sampah berskala besar.

Teknik ketiga adalah bentuk pengaburan paling kasar. Model menolak menjawab untuk masukan
yang dianggap aneh, misalnya yang jaraknya ke pusat data terlalu jauh, lalu mengembalikan
label ``tidak diketahui''. Bagi penyerang ini bukan gradien yang buruk, melainkan tidak
ada gradien: fungsinya fungsi indikator, yang turunannya nol di semua tempat kecuali di
permukaan keputusan. Pertahanan semacam ini juga menolak data baru yang sah, dan data sah
adalah alasan sistem itu dipasang.

Teknik keempat dan kelima adalah keluarga acak. Pada yang keempat, masukan diubah acak
sebelum masuk model: gambar diperbesar, digeser satu piksel, atau sebagian fiturnya
dihapus. Pada yang kelima, yang diacak adalah modelnya: satu ansambel dipilih acak pada
setiap panggilan, atau prediksi diambil dari rata-rata banyak model. Pada keduanya,
gradien satu panggilan adalah gradien satu realisasi, sedangkan yang dihadapi penyerang
adalah sebaran realisasi. Sampel gradien buruk bukan karena bias, melainkan karena
variansi: tanda dari arah yang berubah-ubah bukan arah yang baik.

\index{penyaringan acak}
Teknik keenam menyerang dari dalam. Pelatihan sengaja membuat
gradien kecil di sekitar titik data, misalnya dengan menghukum norma gradien atau
memakai aktivasi yang jenuh, sehingga permukaan kerugian sangat datar di sekitar contoh
dan curam di tempat lain. Penyerang yang melangkah sejauh $\varepsilon$ di daerah datar
tidak ke mana-mana, dan gradiennya memang kecil, tetapi bukan karena modelnya benar.
Teknik ketujuh adalah kerabatnya: masukan dibersihkan lebih dulu oleh prosedur
optimisasi banyak langkah sebelum diklasifikasikan, dan gradien dari pipa itu mudah
meledak.

Gambar~\ref{fig:permukaan} memperlihatkan bentuk permukaan yang lahir dari kuantisasi
pada dua koordinat. Setiap dataran mendatar di sana adalah satu anak tangga. Di dalam
satu dataran, kerugian tidak berubah sama sekali, sehingga penyerang yang melangkah
kecil berhenti di tempat --- tiga panah pendek pada gambar itu. Penyerang yang boleh
melangkah cukup jauh bisa menyeberang ke dataran berikutnya, dan panah putus-putus
menunjukkan penyeberangan itu.

\begin{figure}[!htbp]
\centering
\begin{tikzpicture}
\begin{axis}[
  width=0.80\linewidth, height=7.2cm,
  view={52}{24},
  xlabel={$u$}, ylabel={$v$}, zlabel={$L$},
  domain=-0.35:0.35, y domain=-0.35:0.35,
  samples=27, samples y=27,
  zmin=0, zmax=0.52,
  tick label style={font=\scriptsize},
  label style={font=\small},
  colormap={peta}{color=(layerteal!10) color=(layerteal!60)},
  colorbar=false, axis line style={layergrey}, tick style={layergrey},
]
\addplot3[surf, shader=faceted, faceted color=layergrey!40, line width=0.15pt]
  {0.1*round(10*((x-0.15)^2+(y-0.10)^2))};
\draw[-latex, accent, line width=1pt] (axis cs:0.02,0.02,0) -- (axis cs:0.10,0.06,0);
\draw[-latex, accent, line width=1pt] (axis cs:0.02,0.02,0) -- (axis cs:-0.02,0.08,0);
\draw[-latex, accent, line width=1pt] (axis cs:0.02,0.02,0) -- (axis cs:0.04,-0.05,0);
\node[circle, fill=accent, inner sep=1.5pt] at (axis cs:0.02,0.02,0) {};
\draw[-latex, layerteal, line width=1.2pt, dashed]
  (axis cs:0.02,0.02,0) -- (axis cs:0.30,0.30,0.10);
\node[font=\footnotesize, accent, anchor=north west] at (axis cs:0.09,0.05,0)
  {$\nabla L=0$};
\node[font=\footnotesize, layerteal, anchor=south west] at (axis cs:0.28,0.28,0.12)
  {BPDA};
\end{axis}
\end{tikzpicture}
\caption{Permukaan kerugian yang dibuat datar oleh kuantisasi:
$L(u,v)=\frac{1}{10}\operatorname{round}(10\,s)$ dengan
$s=(u-0{,}15)^2+(v-0{,}10)^2$.
Tiap dataran setinggi $0{,}1$ adalah satu anak tangga: di dalamnya $\nabla L=0$,
jadi tiga panah pendek dari titik merah tidak mengubah ketinggian sama sekali.
Panah putus-putus menyeberang ke anak tangga berikutnya, dan hanya penyerang yang
bersedia melangkah sejauh jarak antar-batas yang sampai ke sana.}
\label{fig:permukaan}
\end{figure}

\begin{tangan}
Bahan: satu penggaris panjang, satu uang logam, satu teman, satu meja.

Ganjal satu ujung penggaris dengan buku sehingga landainya miring, lalu letakkan uang
logam di tengahnya. Tutup matamu dan tahan uang itu dengan satu jari. Temanmu mengetuk
meja pelan-pelan, tidak beraturan, dan tugasmu menebak ke arah mana penggaris menanjak
--- hanya dari getaran yang sampai ke jarimu.

Ulangi sepuluh kali dan catat benar atau salah. Lalu minta temanmu mengetuk lebih keras,
dan ulangi sepuluh kali. Lalu minta ia berhenti mengetuk sama sekali, dan ulangi lagi.

Yang harus kamu lihat: ketukan keras tidak membuat tebakanmu benar. Ketukan yang berhenti
membuat tebakanmu murni tebakan, sekitar separuh benar, tidak lebih baik daripada melempar
koin. Getaran memberitahu bahwa ada sesuatu yang bergerak; ia tidak memberitahu arah.
Itulah bedanya gradien yang buruk dan gradien yang tidak ada.
\end{tangan}

Ketujuh teknik itu menghasilkan tujuh gejala. Lima pertama adalah daftar Athalye dan
rekan-rekan; dua terakhir kami tambahkan dari tipologi yang sama.

\begin{enumerate}
  \item \textbf{Serangan tidak dibatasi.} Laporan mengizinkan gangguan sebesar apa pun,
        lalu menyimpulkan ketahanan dari serangan yang tidak memperhatikan batas.
  \item \textbf{Serangan memakai satu sampel acak.} Pada pertahanan yang mengacak,
        penyerang menarik satu realisasi lalu melaporkan kegagalannya. Yang diukur adalah
        satu lemparan dadu, bukan peluang menang.
  \item \textbf{Serangan hanya satu langkah.} Serangan satu langkah mudah terhenti oleh
        permukaan kerugian yang tidak rata, dan ketahanan yang disimpulkannya hilang
        begitu langkahnya diulang.
  \item \textbf{Menaikkan herotan tidak menaikkan keberhasilan.} Bila anggaran diperbesar
        tetapi keberhasilan tidak bergerak, yang rusak biasanya gradiennya.
  \item \textbf{Gradien dilaporkan nol atau NaN.} Ini gejala paling telanjang: gradien
        yang seluruhnya nol pada model tak konstan menandakan potongan tak
        terdiferensiasi, bukan permukaan yang rata.
  \item \textbf{Akurasi di bawah serangan sama dengan akurasi bersih.} Tambahan kami.
        Kesamaan sampai beberapa angka di belakang koma pada beberapa $\varepsilon$
        berbeda adalah tanda serangan yang tidak bergerak.
  \item \textbf{Tidak ada serangan adaptif.} Tambahan kami. Pertahanan diuji hanya
        dengan serangan yang dirancang sebelum pertahanan itu ada, sehingga belum diuji
        oleh penyerang yang tahu dirinya sedang diuji.
\end{enumerate}

\kamus{gradien dikaburkan}{obfuscated gradient}{gradien yang diberikan kepada penyerang
tidak menunjuk arah yang benar, biasanya karena ada potongan tak terdiferensiasi atau
acak di dalam pipa, bukan karena modelnya kuat}

\kamus{kuantisasi}{quantization}{penggantian setiap nilai dengan kelipatan terdekat dari
suatu jarak tetap; hasilnya fungsi tangga yang turunannya nol di hampir semua tempat}

\kamus{serangan adaptif}{adaptive attack}{serangan yang dirancang dengan mengetahui
pertahanan yang sedang diuji, termasuk mengetahui bahwa gradiennya mungkin rusak}

Daftar gejala itu berguna untuk menolak klaim lemah, tetapi tidak untuk menerima klaim
kuat: pertahanan bisa lolos ketujuhnya dan tetap rapuh, dan sebaliknya bisa memicu satu
gejala dan tetap menaikkan biaya serangan.


\section{Menembus: BPDA, EOT, dan reparameterisasi}
\tingkat{lanjut}

Kalau gradien yang tersedia salah, penyerang punya tiga pilihan. Ganti gradien yang salah
dengan gradien fungsi yang mendekatinya; rata-ratakan gradien atas keacakan yang membuatnya
berubah-ubah; atau pindahkan pekerjaan ke ruang lain, di mana pertahanannya tidak
mengaburkan apa pun. Ketiganya punya nama, dan ketiganya bekerja karena alasan berbeda.

\index{bpda}
Pilihan pertama disebut \emph{backward pass differentiable approximation}, disingkat
BPDA. Idenya sesederhana bunyinya: pada perhitungan maju pakai fungsi yang sebenarnya,
termasuk pembulatan atau penyaringan; pada perhitungan mundur lewati potongan itu
seolah-olah ia identitas. Bila model terbungkus adalah $h(x)=f\bigl(g(x)\bigr)$,
penyerang memakai

\begin{equation}
\widehat{\nabla}_x h(x) \;=\; \nabla f\bigl(g(x)\bigr),
\end{equation}

yaitu gradien $f$ yang dievaluasi di titik setelah dibungkus, bukan gradien dari
komposisi. Untuk pembulatan, ini berarti penyerang berhenti bertanya ``berapa turunan
pembulatan'' dan mulai bertanya ``kalau angkanya sudah dibulatkan, ke arah mana $f$
paling cepat naik''.

Mengapa pendekatan ini masuk akal? Karena kesalahannya dibatasi oleh seberapa jauh $g$
menggeser titik, bukan oleh seberapa aneh bentuk $g$.

\begin{proposisi}
Misalkan $f$ terdiferensiasi kontinu dengan gradien yang Lipschitz, yaitu
$\lVert \nabla f(u)-\nabla f(v)\rVert \le L\lVert u-v\rVert$ untuk semua $u,v$, dan
misalkan $g$ memenuhi $\lVert g(x)-x\rVert\le\delta$. Maka
\begin{equation}
\bigl\lVert \nabla f\bigl(g(x)\bigr)-\nabla f(x)\bigr\rVert \;\le\; L\,\delta .
\end{equation}
\end{proposisi}

\begin{proof}
Tulis $v=g(x)-x$ dan gunakan teorema nilai rata-rata dalam bentuk integral,
\begin{equation}
\nabla f(x+v)-\nabla f(x) = \int_0^1 H_f(x+tv)\,v\,dt,
\end{equation}
dengan $H_f$ matriks Hessian. Karena norma spektral $H_f$ dibatasi $L$,
\begin{equation}
\bigl\lVert \nabla f(x+v)-\nabla f(x)\bigr\rVert
\le \int_0^1 \lVert H_f(x+tv)\rVert\,\lVert v\rVert\,dt
\le L\lVert v\rVert \le L\delta .
\end{equation}
\end{proof}

Proposisi ini menjelaskan kapan BPDA boleh dipercaya. Kalau pembulatan hanya menggeser
masukan sejauh $q/2$ dan $f$ cukup mulus pada skala itu, galat gradiennya kecil dan
langkah serangan tetap menuju arah yang benar. Kalau $q$ besar sehingga penggeseran itu
sebanding dengan skala lekukan $f$, batas $L\delta$ menjadi besar dan tidak ada jaminan.
Proposisi ini tidak mengatakan $\widehat{\nabla}$ adalah gradien $h$; gradien itu tidak
ada. Ia hanya mengatakan bahwa arah yang dipakai dekat dengan arah yang benar untuk $f$
di titik yang sudah dibungkus.

Pilihan kedua adalah \emph{expectation over transformation}, disingkat EOT. Untuk
pertahanan acak, penyerang tidak menghadapi satu fungsi melainkan sebaran fungsi.
Yang masuk akal dimaksimalkan bukan kerugian pada satu realisasi, melainkan harapan
kerugian atas realisasi,

\begin{equation}
\overline{L}(x) \;=\; \mathbb{E}_{T\sim\mathcal{T}}\; L\bigl(f_T(x),y\bigr),
\end{equation}

dengan $T$ menyatakan keacakan yang dipakai pertahanan, misalnya topeng fitur atau
parameter transformasi.

\index{eot}
\begin{lemma}
Bila $\mathcal{T}$ berhingga dan $\nabla_x L\bigl(f_T(x),y\bigr)$ terbatas seragam
dalam $T$, maka
\begin{equation}
\nabla_x \overline{L}(x) \;=\; \mathbb{E}_{T}\;\nabla_x L\bigl(f_T(x),y\bigr).
\end{equation}
Selain itu, rata-rata $m$ sampel bebas,
$\frac{1}{m}\sum_{i=1}^{m}\nabla_x L\bigl(f_{T_i}(x),y\bigr)$, adalah penaksir tak bias
dengan variansi yang mengecil sebanding $1/m$.
\end{lemma}

\begin{proof}
Karena $\mathcal{T}$ berhingga, harapan adalah jumlah berhingga, dan turunan boleh
masuk ke dalam jumlah karena turunan adalah operator linear. Penaksir rata-rata punya
harapan yang sama dengan tiap sukunya, jadi tak bias, dan variansinya adalah variansi
satu suku dibagi $m$ karena sampelnya bebas.
\end{proof}

Lemma ini menjawab keberatan yang sering muncul: kalau pertahanannya acak, bukankah
serangan apa pun hanya kebetulan? Tidak, selama penyerang boleh memanggil model
berkali-kali. Keacakan menaikkan biaya, karena penyerang perlu $m$ panggilan per langkah,
tetapi tidak mengubah arah yang dituju; yang berubah hanya ketelitian taksiran. Ini
contoh pertama dari pola yang berulang sepanjang bab: pertahanan yang menaikkan biaya itu
nyata, sedangkan yang hanya menyulitkan pengukuran itu semu.

Pilihan ketiga adalah reparameterisasi. Alih-alih menyerang $x$, penyerang menyerang
variabel yang menghasilkan $x$ melalui transformasi yang tidak dikaburkan. Kalau
pertahanan membungkus model dengan fungsi $g$ yang punya invers, penyerang bekerja di
ruang $z$ dengan $x=g^{-1}(z)$ dan mengoptimalkan $z$. Ini berguna ketika $g$ adalah
rotasi atau transformasi linier lain yang tidak menghilangkan informasi, dan tidak
berguna ketika $g$ menghilangkan informasi: tidak ada $z$ yang bisa memulihkan angka
yang sudah dibuang.

Dalam praktik, ketiga pilihan itu dipakai bersama, ditambah dua kebiasaan. Pertama: mulai
dari beberapa titik awal acak di dalam bola $\varepsilon$, karena permukaan kerugian yang
tidak cembung membuat satu titik awal mudah terjebak. Kedua: jangan berhenti pada satu
nilai $\varepsilon$. Ketahanan yang diukur pada satu anggaran adalah ketahanan yang tidak
diukur.

\begin{resep}
Empat langkah menguji klaim pertahanan, dalam urutan ini.
\begin{enumerate}
  \item \textbf{Tulis anggaran lebih dulu.} Tetapkan $\varepsilon$ atau batas byte,
        jumlah kueri, dan waktu. Angka ini tidak boleh berubah setelah hasil pertama
        keluar.
  \item \textbf{Serang lewat transfer, tanpa kueri.} Latih model pengganti pada data
        yang kamu punya, buat contoh adversarial di sana, lalu kirim ke model korban.
        Langkah ini mengukur apakah pertahanan itu hanya menghalangi satu model.
  \item \textbf{Belanjakan sisa kueri secara adaptif.} Baru sekarang pakai BPDA atau
        EOT pada korban, dan catat kueri per langkah, bukan hanya hasil akhirnya.
  \item \textbf{Laporkan keduanya.} Keberhasilan transfer, keberhasilan adaptif, dan
        biaya masing-masing, selalu bersama akurasi bersih.
\end{enumerate}
Peringatan urutan: langkah 2 harus selesai sebelum langkah 3. Kalau kamu mengukur
transfer setelah melihat hasil adaptif, kamu sudah tahu contoh mana yang berhasil,
dan angka transfernya bocor ke depan.
\end{resep}

\kamus{bpda}{backward pass differentiable approximation}{teknik mengganti gradien
potongan yang tidak terdiferensiasi dengan gradien fungsi di titik sesudah potongan itu,
seolah potongan tersebut identitas}

\kamus{eot}{expectation over transformation}{serangan yang memaksimalkan harapan
kerugian atas keacakan pertahanan, biasanya dengan merata-ratakan gradien dari beberapa
realisasi}

\begin{untukdosen}
Bagian ini biasanya diajarkan dalam urutan yang salah, yaitu definisi BPDA lebih dulu.
Urutan yang lebih baik adalah sebaliknya: tunjukkan model yang akurasinya tidak bergerak
sedikit pun ketika diserang, minta mahasiswa menebak mengapa, lalu tunjukkan gradiennya
nol. Hampir semua akan menebak ``modelnya kuat'', dan keheningan sesudahnya adalah momen
belajar tersendiri.

Dua kesalahan perlu diantisipasi. Pertama, menyimpulkan bahwa BPDA ``menghapus''
pertahanan; yang benar, BPDA mengembalikan akses ke arah yang kira-kira benar, dan
kegunaannya bergantung pada besarnya $\delta$ relatif terhadap skala lekukan model.
Kedua, menyimpulkan bahwa pertahanan acak tidak berguna karena EOT ada; yang benar,
pertahanan acak menaikkan biaya per langkah secara linear terhadap $m$, sehingga
pertahanan itu mengubah anggaran penyerang, bukan arah serangan. Latihan di akhir bab
memaksa mahasiswa menghitung biaya itu, bukan menyebutnya.
\end{untukdosen}

\section{Contoh hitung: satu titik pada model mainan}
\tingkat{menengah}

\begin{untukanak}
Seorang penjaga toko punya kebiasaan aneh: setiap harga dibulatkan ke rupiah terdekat
sebelum dicatat. Kalau ada pembeli yang diam-diam menggeser harga sebesar satu sen,
catatan penjaga tidak berubah sama sekali. Pembeli itu bisa menyimpulkan bahwa tidak
ada gunanya menggeser harga, lalu berhenti mencoba. Kesimpulannya salah, tetapi
alasannya masuk akal: yang ia ukur memang tidak bergerak.

Sekarang bayangkan pembeli lain. Ia tidak bertanya apakah catatan berubah karena
geseran satu sen; ia bertanya ke arah mana catatan paling cepat naik, lalu menggeser
harga cukup jauh sehingga pembulatannya berubah. Catatan penjaga pun berubah, dan
pembeli itu berhasil tanpa melanggar aturan apa pun. Yang membedakan keduanya bukan
kepandaian, melainkan pertanyaan yang diajukan.

Model yang memakai pembulatan bekerja seperti penjaga itu. Program yang hanya mencoba
geseran kecil akan berhenti dan melaporkan bahwa modelnya aman; program yang mengerti
pembulatan akan melangkah lebih jauh. Jadi ketika membaca laporan bahwa sebuah model
tidak bisa dikalahkan, tanyakan dulu: penyerangnya berhenti karena kalah, atau karena ia
hanya mencoba satu cara?
\end{untukanak}

Sebelum percaya pada tabel, kita hitung satu titik dengan tangan. Modelnya dibuat
sesederhana mungkin: delapan fitur dengan bobot semuanya satu, jadi skornya adalah
jumlah kedelapan fitur. Datanya berasal dari dua sebaran normal berpusat $-0{,}5$ dan
$+0{,}5$, dengan akurasi bersih sekitar $92$ persen. Pertahanannya adalah pembulatan ke
kelipatan $q=0{,}05$ pada setiap fitur, ditambah penyaringan acak berpeluang hapus $r$
yang mengalikan sisanya dengan $1/(1-r)$. Anggaran penyerang $\varepsilon=0{,}25$ pada
norma $\ell_\infty$.

Ambil satu contoh yang benar berlabel $+1$,

\begin{equation}
x_0 = (0{,}30;\ 0{,}20;\ 0{,}15;\ 0{,}10;\ 0{,}05;\ 0{,}00;\ -0{,}10;\ -0{,}08).
\end{equation}

Jumlahnya $0{,}62$, jadi model menjawab $+1$ dengan margin yang tipis. Sesudah
pembulatan, vektor yang masuk ke pengklasifikasi adalah
$(0{,}30;0{,}20;0{,}15;0{,}10;0{,}05;0{,}00;-0{,}10;-0{,}10)$ dan skornya $0{,}60$.
Perhatikan bahwa satu koordinat bergeser dari $-0{,}08$ ke $-0{,}10$: itulah kerja
pembulatan.

Sekarang bandingkan dua penyerang. Penyerang naif menghitung turunan fungsi yang
sebenarnya, yaitu turunan dari komposisi yang memuat pembulatan. Karena pembulatan
datar di sekitar $x_0$, turunannya nol di kedelapan koordinat, sehingga langkahnya
$x_0 + \varepsilon\cdot 0 = x_0$. Skornya tetap $0{,}60$ dan labelnya tetap $+1$.
Laporan penyerang itu akan berbunyi: serangan gagal, akurasi bertahan seratus persen
dari akurasi bersih.

Penyerang kedua memakai BPDA. Ia tidak bertanya tentang turunan pembulatan; ia
mengevaluasi gradien pengklasifikasi di titik yang sudah dibulatkan, yaitu vektor
$(1,1,\dots,1)$, lalu melangkah berlawanan arah label sejauh $\varepsilon$ pada setiap
koordinat,

\begin{equation}
x_1 = x_0 - 0{,}25\cdot(1,1,\dots,1)
= (0{,}05;\ -0{,}05;\ -0{,}10;\ -0{,}15;\ -0{,}20;\ -0{,}25;\ -0{,}35;\ -0{,}33).
\end{equation}

Sesudah pembulatan, vektor itu menjadi
\[
(0{,}05;\ -0{,}05;\ -0{,}10;\ -0{,}15;\ -0{,}20;\ -0{,}25;\ -0{,}35;\ -0{,}35)
\]
dengan skor $-1{,}40$. Labelnya berbalik menjadi $-1$. Satu langkah, biaya satu kali
perhitungan gradien, dan model yang tadi dilaporkan kebal ternyata salah dengan margin
yang tidak tipis.

Contoh ini menunjukkan sesuatu yang lebih tajam daripada ``BPDA bekerja''. Penyerang
kedua tidak memakai informasi yang tidak dimiliki penyerang pertama: keduanya tahu
arsitektur, bobot, dan $q$. Yang berbeda hanya satu keputusan, yaitu menghormati potongan
tak terdiferensiasi atau melewatinya. Jadi ``ketahanan'' yang dilaporkan penyerang pertama
bukan properti model, melainkan properti notasi penyerang itu sendiri.

\begin{contoh}
Pada model yang sama, dengan $q=0{,}05$, $r=0{,}5$, dan seed tetap, gradien beda
hingga dengan langkah $h=10^{-5}$ bernilai nol pada $15\,998$ dari $16\,000$ entri;
dua entri sisanya tepat $q/h=5\,000$, karena titiknya kebetulan berada di batas anak
tangga. Norma gradien beda hingga yang diambil nilai tengahnya nol, sedangkan norma
gradien BPDA untuk satu topeng yang menyisakan empat fitur adalah $4$. Tetapi arah satu
topeng masih jauh dari arah bobot, dengan kosinus antara $0{,}35$ dan $0{,}94$ atas
$200$ topeng. Rata-rata gradien atas $200$ topeng memberi vektor bernorma $2{,}738$ dan
kosinus $0{,}9982$ terhadap vektor bobot.
\end{contoh}

Angka itu menggambarkan tiga hal. Pertama, gradien yang hancur bukan hanya kecil; ia bisa
juga besar dan salah, dan dua entri yang melonjak akan mendominasi setiap langkah
sehingga serangan naif gagal dengan cara yang tampak acak. Kedua, BPDA memulihkan arah
dengan hampir sempurna: kosinus $0{,}9982$ berarti sudutnya sekitar tiga derajat. Ketiga,
norma gradien BPDA untuk satu topeng, yaitu $4$, berbeda dari norma harapannya, $2{,}738$;
perbedaan skala itu tidak merusak serangan yang memakai tanda, tetapi merusak yang memakai
besaran gradien, dan itu sebabnya EOT tetap perlu.

\kertas{Hitung ulang contoh di atas tanpa melihat teks. Jumlahkan $x_0$, bulatkan ke
kelipatan $0{,}05$, lalu jumlahkan lagi; cocokkan bahwa hasilnya $0{,}60$. Kurangi
setiap koordinat dengan $0{,}25$, bulatkan lagi, jumlahkan, dan cocokkan bahwa hasilnya
$-1{,}40$. Sesudah itu jawab: berapa nilai $q$ terkecil yang membuat serangan BPDA satu
langkah ini gagal membalik label, dan mengapa nilai itu bergantung pada margin awal
$0{,}60$ dan bukan pada $\varepsilon$ saja?}

\section{Lab komputasi: akurasi di tiap tingkat pengaburan}
\tingkat{lanjut}

\begin{figure}[!htb]
\centering
\begin{tikzpicture}
\begin{axis}[
  width=0.86\linewidth, height=5.6cm,
  xlabel={anggaran $\varepsilon$ pada norma $\ell_\infty$},
  ylabel={keberhasilan serangan (\%)},
  xmin=0.03, xmax=0.37, ymin=-0.6, ymax=30,
  xtick={0.05,0.10,0.15,0.20,0.25,0.30,0.35},
  xticklabels={0{,}05,0{,}10,0{,}15,0{,}20,0{,}25,0{,}30,0{,}35},
  ytick={0,5,10,15,20,25,30},
  tick label style={font=\footnotesize},
  label style={font=\small},
  legend style={font=\scriptsize, at={(0.02,0.98)}, anchor=north west,
                draw=layergrey!35, fill=white, fill opacity=0.92, text opacity=1},
  grid=major, grid style={layergrey!15},
  axis line style={layergrey},
]
\addplot[color=layerindigo, mark=*, mark size=1.5pt, line width=1pt]
  coordinates {(0.05,2.38) (0.10,5.27) (0.15,8.66) (0.20,12.84)
               (0.25,17.40) (0.30,22.32) (0.35,27.71)};
\addlegendentry{tanpa pembulatan, penyerang BPDA}
\addplot[color=accent, mark=square*, mark size=1.4pt, line width=1pt, dashed]
  coordinates {(0.05,2.34) (0.10,5.33) (0.15,8.72) (0.20,12.92)
               (0.25,17.48) (0.30,22.46) (0.35,27.79)};
\addlegendentry{pembulatan $q=0{,}05$, penyerang BPDA}
\addplot[color=layergrey, mark=triangle*, mark size=1.6pt, line width=0.9pt,
         dotted]
  coordinates {(0.05,0.00) (0.10,0.01) (0.15,0.01) (0.20,0.01)
               (0.25,0.01) (0.30,0.02) (0.35,0.02)};
\addlegendentry{pembulatan $q=0{,}05$, penyerang naif}
\end{axis}
\end{tikzpicture}
\caption{Keberhasilan serangan terhadap $\varepsilon$ pada model mainan delapan fitur,
$2\,000$ contoh, rata-rata lima seed. Pada $\varepsilon=0{,}25$ ketiga kurva memberi
$17{,}40$, $17{,}48$, dan $0{,}01$ persen; pada $\varepsilon=0{,}35$ memberi
$27{,}71$, $27{,}79$, dan $0{,}02$ persen. Dua kurva pertama berimpit di dalam
simpangan antar-seed, yang besarnya $0{,}35$ sampai $1{,}08$ poin persen.}
\label{fig:kurva-epsilon}
\end{figure}

Sekarang kita ulang percobaan itu pada dua ribu contoh, lima seed, dan beberapa tingkat
pengaburan, hanya dengan \texttt{numpy}. Yang perlu diperhatikan bukan kode itu sendiri,
melainkan satu keputusan di dalamnya: fungsi \texttt{fgsm} menghitung gradien dengan beda
hingga pada fungsi yang sudah dibungkus, sedangkan \texttt{pgd} memakai gradien
pengklasifikasi di titik sesudah pembungkusan. Hasilnya, disapu pada tujuh nilai
$\varepsilon$, ada di Gambar~\ref{fig:kurva-epsilon}.

\begin{lab}
\begin{verbatim}
import numpy as np

D, N, EPS = 8, 2000, 0.25
W = np.ones(D)

def data(rng, n, mu=0.5):
    y = 2 * rng.integers(0, 2, n) - 1
    return rng.normal(mu * y[:, None], 1.0, (n, D)), y

def quant(x, q):
    return x if q == 0 else np.round(x / q) * q

def mask(shape, r, rng):
    if r == 0:
        return np.ones(shape)
    return (rng.random(shape) >= r) / (1 - r)

def logit(x, q, m):
    return (quant(x, q) * m) @ W

def fgsm(X, y, q, r, rng):
    m, h = mask(X.shape, r, rng), 1e-5
    G = np.zeros_like(X)
    for j in range(D):
        e = np.zeros(D)
        e[j] = h
        G[:, j] = (logit(X + e, q, m) - logit(X - e, q, m)) / (2 * h)
    return X - EPS * y[:, None] * np.sign(G)

def pgd(X, y, q, r, rng, m_eot=1, steps=40):
    Xa, a = X.copy(), EPS / 8
    for _ in range(steps):
        G = sum(W * mask(X.shape, r, rng) for _ in range(m_eot))
        Xa = np.clip(Xa - a * y[:, None] * np.sign(G / m_eot),
                     X - EPS, X + EPS)
    return Xa

def run(seed, q, r):
    rng = np.random.default_rng(seed)
    X, y = data(rng, N)
    def pred(Z):
        return np.sign(logit(Z, q, mask(Z.shape, r, rng)))
    ok = pred(X) == y
    def succ(Z):
        return 100 * np.mean(pred(Z)[ok] != y[ok])
    return (100 * ok.mean(), succ(fgsm(X, y, q, r, rng)),
            succ(pgd(X, y, q, r, rng)), succ(pgd(X, y, q, r, rng, 20)))

print("  q     r   bersih   naif   bpda1   eot")
for q, r in ((0.00, 0.0), (0.01, 0.0), (0.05, 0.0),
             (0.05, 0.5), (0.20, 0.0)):
    M = np.array([run(s, q, r) for s in (1, 2, 3, 4, 5)]).mean(0)
    print("%5.2f %5.2f  %6.2f  %5.2f  %5.2f  %5.2f"
          % (q, r, M[0], M[1], M[2], M[3]))
\end{verbatim}

Keluaran program itu, apa adanya:

\begin{verbatim}
  q     r   bersih   naif   bpda1   eot
 0.00  0.00   92.02  17.40  17.40  17.40
 0.01  0.00   92.00   0.02  17.40  17.40
 0.05  0.00   91.97   0.01  17.48  17.48
 0.05  0.50   82.87  12.16  25.48  25.67
 0.20  0.00   91.77   0.00  17.65  17.65
\end{verbatim}
\end{lab}

Bacalah tabel itu baris demi baris. Baris pertama adalah garis dasar tanpa pengaburan:
akurasi bersih $92{,}02$ persen, keberhasilan serangan satu langkah $17{,}40$ persen.
Angka itu kecil karena $\varepsilon=0{,}25$ memang kecil dibanding sebaran datanya,
bukan karena serangannya lemah.

Baris kedua adalah inti bab ini. Dengan $q=0{,}01$, akurasi bersih hampir tidak berubah,
$92{,}00$ persen, tetapi keberhasilan naif jatuh dari $17{,}40$ menjadi $0{,}02$ persen,
sekitar $870$ kali lebih kecil. Penyerang yang memakai BPDA mendapat $17{,}40$ persen,
persis garis dasar. Pertahanan itu tidak mengubah apa pun yang bisa diamati penyerang
yang menghitung dengan benar; ia hanya mengubah apa yang dilihat penyerang yang salah
menghitung.

Baris ketiga mengulang hal yang sama pada $q=0{,}05$: naif $0{,}01$ persen, BPDA
$17{,}48$ persen, akurasi bersih $91{,}97$ persen. Menaikkan $q$ lima kali lipat tidak
mengubah keberhasilan BPDA, tetapi membuat laporan naif terlihat lebih baik. Itulah
gejala keempat dari daftar di bagian dua, dalam bentuk angka.

Baris keempat menambahkan penyaringan acak $r=0{,}5$, dan dua hal terjadi sekaligus.
Akurasi bersih turun dari $91{,}97$ ke $82{,}87$ persen, karena membuang separuh fitur dan
mengalikan sisanya dengan dua menaikkan variansi skor. Keberhasilan BPDA naik ke $25{,}48$
persen, di atas garis dasar, bukan karena pertahanannya lebih buruk melainkan karena contoh
yang masih benar setelah pengaburan punya margin lebih kecil.

Baris kelima menunjukkan bahwa kuantisasi kasar tidak menyelamatkan model. Dengan
$q=0{,}20$, sebanding dengan $\varepsilon$, keberhasilan BPDA $17{,}65$ persen, sedikit di
atas garis dasar: titik yang berada dekat batas anak tangga hanya perlu dorongan kecil untuk
menyeberang.

Variansi antar seed kecil tetapi tidak nol: pada $q=0{,}05$ dengan $r=0$ simpangan bakunya
sekitar $0{,}6$ poin persen, pada baris dengan penyaringan acak sekitar $1{,}4$ poin. Galat
baku satu proporsi pada $2\,000$ contoh dan tingkat $0{,}17$ kira-kira $0{,}8$ poin, jadi
kedua angka itu sejalan dengan pengambilan sampel biasa.

Ada empat hal yang tidak diukur percobaan ini. Pertama, modelnya linier, sehingga gradien
BPDA hampir sempurna dan kosinus $0{,}9982$ terlalu bagus untuk model nyata. Kedua,
penyerang mengetahui $q$ dan $r$. Ketiga, hanya satu nilai $\varepsilon$ yang diuji.
Keempat, tidak ada pencarian hyperparameter serangan. Semua itu membuat angka di atas
menjadi batas bawah keberhasilan serangan, bukan batas atas.


Kekurangan ketiga itu bisa ditutup, dan hasilnya ada di Gambar~\ref{fig:kurva-epsilon}.
Program yang sama dijalankan pada tujuh nilai $\varepsilon$ dari $0{,}05$ sampai
$0{,}35$, lima seed, tanpa penyaringan acak. Kurva model tanpa pembulatan dan kurva
model dengan pembulatan $q=0{,}05$ hampir berimpit: pada $\varepsilon=0{,}25$ keduanya
memberi $17{,}40$ dan $17{,}48$ persen, dan pada $\varepsilon=0{,}35$ keduanya memberi
$27{,}71$ dan $27{,}79$ persen. Kurva penyerang naif tetap di bawah $0{,}02$ persen
sepanjang rentang itu. Pertahanan ini karena itu tidak menggeser kurva dan tidak
mematahkannya; ia hanya menambah kurva ketiga yang datar.


\section{Studi kasus: NIDS, dan mengapa evasinya tetap tidak praktis}
\tingkat{lanjut}

Sampai di sini kita berurusan dengan vektor fitur yang boleh diubah sesuka hati. Sistem
pengesanan pencerobohan jaringan, atau NIDS, adalah dunia yang berbeda.
\index{nids}
Sebuah NIDS membaca aliran paket, mengekstrak ciri seperti durasi aliran,
jumlah paket, jumlah byte, bendera TCP, dan waktu antar-kedatangan, lalu memutuskan
apakah aliran itu serangan. Keputusan harus diambil cepat, sering dalam milidetik, untuk
ribuan aliran per detik.

Gambar~\ref{fig:rantai} menaruh rantai itu pada satu baris. Yang perlu diperhatikan
bukan kotaknya, melainkan tanda merah di antara rekonstruksi aliran dan ekstraksi ciri.
Di situlah penyerang bekerja: ia tidak menyentuh model, ia menyentuh paket, dan paket
itulah yang menentukan nilai ciri yang masuk ke model. Semua yang ada di kanan tanda itu
hanya melihat akibat; semua yang ada di kiri adalah sebab yang bisa ia kendalikan, dengan
batasan protokol yang dibahas sesudah ini.

\begin{figure}[!htbp]
\centering
\begin{tikzpicture}[font=\footnotesize,
  tahap/.style={draw=layergrey!70, fill=layergrey!6, rounded corners=1.5pt,
    align=center, inner sep=3pt, minimum height=1.1cm, text width=1.95cm},
  panah/.style={-latex, layergrey, line width=0.8pt}]
\node[tahap] (n1) at (0,0) {paket\\di kabel};
\node[tahap] (n2) at (2.50,0) {rekonstruksi\\aliran};
\node[tahap] (n3) at (5.00,0) {ekstraksi\\ciri};
\node[tahap] (n4) at (7.50,0) {penyetelan\\skala};
\node[tahap] (n5) at (10.00,0) {model\\NIDS};
\node[tahap] (n6) at (12.50,0) {ambang\\keputusan};
\draw[panah] (n1) -- (n2);
\draw[panah] (n2) -- (n3);
\draw[panah] (n3) -- (n4);
\draw[panah] (n4) -- (n5);
\draw[panah] (n5) -- (n6);
\node[circle, fill=accent, inner sep=2pt] (sisip) at (3.75,0) {};
\draw[accent, line width=0.8pt] (sisip) -- ++(0,0.62);
\node[accent, anchor=south, font=\scriptsize] at (3.75,0.66)
  {titik sisipan: waktu dan ukuran paket};
\node[layergrey, anchor=north, font=\scriptsize, align=center] at (12.5,-0.62)
  {peringatan atau blokir};
\draw[-latex, layergrey, line width=0.6pt] (n6.south) -- (12.5,-0.5);
\end{tikzpicture}
\caption{Rantai pemrosesan NIDS, dari kabel sampai keputusan. Tanda merah adalah satu
titik sisipan yang bisa dipakai penyerang: ciri seperti jumlah paket, jumlah byte, dan
waktu antar-kedatangan seluruhnya ditentukan di kiri tanda itu, sedangkan model di
kanannya hanya membaca akibatnya. Pertahanan yang menambahkan pengaburan di kanan tanda
tidak menyentuh sebab, dan itulah batas yang dibahas di bagian ini.}
\label{fig:rantai}
\end{figure}

\kamus{nids}{network intrusion detection system}{sistem yang memantau lalu lintas
jaringan dan memutuskan apakah suatu aliran paket merupakan serangan}

\kamus{ruang masalah}{problem space}{ruang objek nyata yang boleh diubah penyerang,
misalnya paket dan waktu, berbeda dari ruang fitur tempat model bekerja}

Kesulitan pertama adalah jarak antara dua ruang. Serangan biasanya dirumuskan di ruang
fitur, sedangkan penyerang mengirim paket: setiap perubahan fitur harus diterjemahkan
menjadi perubahan paket yang tetap sah menurut protokol dan tetap menjalankan fungsi
jahatnya, dan terjemahan itu tidak selalu ada. Beberapa fitur bahkan tidak bisa diubah
pengirim, karena jumlah paket dan byte arah balik ditentukan sistem penerima.

Kesulitan kedua adalah korelasi antar fitur. Durasi berkaitan dengan jumlah paket dan
waktu antar-kedatangan, dan bendera TCP berkaitan dengan protokol transpor, sehingga
penyerang tidak bisa menaikkan durasi sambil menurunkan kedua waktu antar-kedatangan;
fitur diskret seperti nomor protokol dan panjang paket menambah batas.

Kesulitan ketiga adalah waktu. Serangan yang membutuhkan sepuluh ribu kueri berarti lalu
lintas nyata yang lewat dan tercatat, dan jejak adalah hal pertama yang dilihat analis.
Biaya nyata adalah pertahanan nyata; ketiadaan gradien bukan.

Angka dari literatur membuat jarak itu konkret. Pada data CSE-CIC-IDS2018 dengan lima
puluh fitur tersaring, tiga model dilatih dan diuji pada dua puluh persen data, lalu
diserang dengan FGSM, PGD, dan BIM. Sebelum serangan ketiganya nyaris sempurna: akurasi
$0{,}997$ untuk ANN, $0{,}998$ untuk SVM, dan $0{,}997$ untuk CNN. Pada hari serangan,
dengan pengetahuan penuh, FGSM menjatuhkan ANN ke $0{,}756$ dengan F1 $0{,}677$, CNN ke
$0{,}842$ dengan F1 $0{,}863$, dan SVM ke $0{,}150$ dengan F1 $0{,}000$. PGD sedikit
lebih kuat pada CNN, $0{,}584$ dengan F1 $0{,}689$, dan BIM memberi $0{,}688$ dengan F1
$0{,}751$.

Angka itu terlihat seperti kemenangan penyerang, dan pada model statis memang demikian;
tetapi percobaan yang sama melanjutkan dengan model yang dilatih ulang setiap hari.
Dengan pengetahuan tentang model berumur satu hari, SVM kembali ke $0{,}998$, ANN ke
$0{,}938$, dan CNN ke $0{,}957$ dengan F1 $0{,}959$, dan pada hari kedua angkanya antara
$0{,}966$ dan $0{,}998$. Pelatihan ulang itu tidak memakai contoh adversarial, jadi
pemulihannya bukan hasil pertahanan adversarial melainkan akibat fakta bahwa gradien yang
dihitung pada model kemarin tidak lagi menunjuk arah yang benar.

Inilah versi NIDS dari seluruh bab: pengaburan gradien yang disengaja rapuh, karena
penyerang bisa menembusnya dengan BPDA, sedangkan pengaburan yang datang dari perubahan
model lebih tahan lama, sebab penyerang harus memperbarui pengetahuannya, dan memperbarui
pengetahuan berarti memperoleh akses, dan akses meninggalkan jejak. Bedanya bukan pada
teknik, melainkan pada siapa yang membayar biaya.

\index{transferabilitas}
Penulis tinjauan itu menyusun pohon serangan dan taksonomi masalah kepraktisan, lalu
menyimpulkan bahwa hanya tiga jalur yang layak. Jalur pertama adalah gangguan buta:
menambahkan byte acak pada paket UDP, atau pada paket TCP dengan bendera PSH, sambil
menghitung ulang checksum supaya protokolnya sah. Jalur ini tidak memerlukan pengetahuan
model dan tidak memerlukan kueri, jadi ia kebal terhadap pelatihan ulang, tetapi ia tidak
punya arah. Jalur kedua adalah kueri lewat saluran samping, yang menuntut akses jaringan
berkelanjutan dan bergantung pada indikator yang berubah menurut beban jaringan dan
perangkat keras. Jalur ketiga adalah model pengganti yang dilatih dari data tangkapan
botnet sendiri, yang bekerja lewat transferabilitas tetapi menuntut data latih yang masih
baru.

Ketiga jalur itu nyata, dan penulisnya menekankan bahwa jarak antara penelitian dan
praktik lebar serta bahwa asumsi model ancaman dalam banyak makalah terlalu murah hati
kepada penyerang.

\index{evasi}
Kerja penulis sendiri, yang akan kami sebut catatan kerja dan bukan hasil yang sudah
diterbitkan, bergerak di sisi lain persoalan yang sama. Pada satu kumpulan data aliran
jaringan, sebuah model mencapai F1 $0{,}967037$ dari $46\,558$ atas $48\,145$ contoh uji.
Pertanyaannya bukan ``apakah serangan berhasil'', melainkan ``berapa fitur yang harus
diubah agar sebuah contoh berpindah sisi batas keputusan sambil tetap sah menurut
protokol''. Dengan program linier campuran bilangan bulat, jawabannya bergantung tajam
pada jumlah fitur yang boleh berubah: pada tujuh fitur, tidak satu pun dari lima puluh
contoh bisa diseberangkan; pada dua puluh fitur, tiga puluh tiga dari lima puluh bisa.
Rasio jarak ke batas antara contoh yang berhasil dan contoh kendali sekitar $3{,}28$
untuk model linier dan $2{,}13$ untuk jaringan saraf.

\index{aturan tiga}
Dua angka lagi menunjukkan kehati-hatian. Dari dua puluh ribu kandidat yang diuji untuk
konsistensi protokol tidak ada yang lolos; dengan aturan tiga, batas atas $95$ persen
untuk laju itu sekitar $0{,}019$ persen. Pada model dengan margin keselamatan yang
diperbesar, syaratnya berbentuk $\varepsilon+40\delta$ dengan
$\delta=1{,}245\times10^{-7}$, sehingga anggaran efektifnya $0{,}150004981$ dan bukan
$0{,}15$.

Batas catatan kerja itu harus disebut bersama hasilnya. Ia memakai satu kumpulan data,
lima puluh contoh uji, dan dua keluarga model. Ia memverifikasi ketiadaan serangan yang
memenuhi syarat, bukan menemukan serangan. Angka $0$ dari $20\,000$ adalah pernyataan
tentang generator kandidat yang dipakai, bukan tentang semua paket yang mungkin, dan
hasilnya tidak mengatakan apa pun tentang NIDS dinamis.

Sekarang kesimpulan yang harus diucapkan jelas. Model mainan di bagian lima tidak
membuktikan apa pun tentang NIDS nyata; ia menunjukkan bahwa pembulatan yang menghapus
gradien menghasilkan laporan yang menyesatkan, dan bahwa mekanisme itu bisa ditembus
dengan biaya wajar. Ia tidak mengukur biaya serangan di ruang paket. Yang bisa dibawa
dari bab ini adalah kebiasaan: setiap kali sebuah pertahanan melaporkan keberhasilan nol,
tanyakan gradiennya, lalu tanyakan siapa yang membayar biaya.

\begin{tantangan}
Sebuah makalah melaporkan pertahanan yang menurunkan keberhasilan serangan dari $96$
persen ke $3$ persen, dengan akurasi bersih yang tidak berubah sama sekali. Penyerang
mereka memakai FGSM satu langkah, $\varepsilon=8/255$, gradien eksak, satu titik awal,
dan satu seed. Pertanyaannya: apakah ketahanan itu nyata, dan percobaan termurah apa
yang paling cepat membuktikan bahwa angka $3$ persen itu artefak?

Petunjuk pertama: perhatikan apa yang tidak dibayar. Pertahanan yang benar-benar
menaikkan biaya hampir selalu membayar sesuatu, entah akurasi, entah latensi, entah
jumlah parameter.

Petunjuk kedua: kalau gradien dihitung melalui potongan tak terdiferensiasi, angka
$3$ persen itu bukan keberhasilan serangan, melainkan sisa contoh yang kebetulan lahir
dekat batas keputusan. Ulangi dengan PGD empat puluh langkah, BPDA, dan tiga titik awal;
lalu sapu $\varepsilon$ dan laporkan kurvanya.

Kalau mentok, kembali ke bagian tiga.
\end{tantangan}

\begin{untukpeneliti}
Tiga hal di bagian ini sering salah dikutip. Pertama, kategori gradien hancur, acak, dan
hilang adalah taksonomi dari makalah obfuscated gradients; daftar tujuh teknik di bagian
dua adalah dekomposisi pedagogis kami. Di makalah itu ada tiga kategori, sembilan
pertahanan yang diuji, tujuh di antaranya bergantung pada gradien yang dikaburkan, dan
enam yang ditembus sepenuhnya ditambah satu yang ditembus sebagian.

Kedua, BPDA bukan penyelesaian umum. Batas galat $L\delta$ hanya berlaku bila $f$ mulus
dan $\delta$ kecil terhadap skala lekukan; pada pembulatan dengan $q$ sebanding dengan
$\varepsilon$ jaminan itu kosong. Sebaliknya, EOT bukan obat untuk serangan yang salah
arah: merata-ratakan arah yang salah tetap menghasilkan arah yang salah.

Ketiga, hasil NIDS yang kami kutip berasal dari percobaan ruang fitur dengan pengetahuan
penuh dan pelatihan ulang harian, bukan dari serangan ruang paket yang dijalankan. Ia
mendukung klaim sempit: pelatihan ulang memulihkan akurasi terhadap serangan yang memakai
gradien model lama. Ia tidak mendukung klaim bahwa NIDS kebal terhadap penyerang yang
boleh mengubah paket.
\end{untukpeneliti}

\section{Daftar periksa sebelum menulis bahwa serangan berhasil}
\tingkat{dasar}

Bagian ini adalah ringkasan yang bisa dipakai langsung. Setiap butir adalah pertanyaan
yang harus bisa dijawab sebelum sebuah klaim tentang serangan atau pertahanan masuk ke
abstrak.

\begin{enumerate}
  \item \textbf{Model ancaman ditulis lengkap.} Apa yang diketahui penyerang: arsitektur,
        bobot, data latih, probabilitas, atau hanya label? Berapa kueri yang boleh ia
        lakukan, dan berapa lama?
  \item \textbf{Serangan dibatasi pada anggaran yang masuk akal.} Sebutkan $\varepsilon$
        dan normanya, atau batas byte dan paket bila bekerja di ruang masalah. Serangan
        tanpa batas mengukur sesuatu yang tidak dihadapi siapa pun.
  \item \textbf{Serangan bersifat adaptif.}
        \index{serangan adaptif}
        Penyerang harus mengetahui pertahanan yang diuji, termasuk kemungkinan bahwa
        gradiennya rusak. Pertahanan yang hanya diuji oleh serangan dari makalah
        sebelumnya belum diuji.
  \item \textbf{Gradien diperiksa sebelum dipercaya.} Laporkan norma gradien, bagian
        entri yang nol, dan korelasinya dengan arah yang benar. Gradien nol pada model
        tak konstan adalah temuan, bukan hasil.
  \item \textbf{Keberhasilan dilaporkan bersama akurasi bersih.} Pertahanan yang
        menurunkan akurasi bersih dari $92$ ke $60$ persen tidak boleh dibandingkan
        langsung dengan yang mempertahankan $92$ persen.
  \item \textbf{Beberapa seed, dengan rata-rata dan simpangan.} Satu seed adalah satu
        lemparan dadu; selisih dua baris yang lebih kecil daripada simpangannya bukan
        selisih.
  \item \textbf{Kurva, bukan satu titik.} Sapu nilai $\varepsilon$, jumlah langkah,
        jumlah sampel EOT, dan jumlah kueri. Pertahanan yang baik menggeser kurva;
        yang semu hanya menggeser satu titik.
  \item \textbf{Anggaran komputasi penyerang disebut:} berapa kali model dipanggil,
        berapa jam, berapa GPU? Serangan yang berhasil setelah sejuta kueri bukan
        serangan yang sama dengan yang berhasil setelah seratus.
  \item \textbf{Kode dan bobot tersedia.} Klaim yang tidak bisa diulang tidak bisa
        diperiksa. Bila pelepasan tidak mungkin, prosedur dan hyperparameter harus cukup
        rinci untuk direkonstruksi.
  \item \textbf{Hasil negatif dilaporkan sebagai hasil negatif.} Pertahanan yang gagal
        adalah informasi; laporan yang berhasil karena serangannya dihentikan sebelum
        waktunya adalah kesalahan yang mahal.
\end{enumerate}

\kamus{aturan tiga}{rule of three}{bila tidak ada kejadian dalam $n$ percobaan bebas,
batas atas $95$ persen untuk peluang kejadian itu kira-kira $3/n$}

\begin{ringkasan}
\begin{itemize}
  \item Gradien yang dikaburkan membuat penyerang buta, bukan membuat model kuat.
        Ketujuh teknik itu menghasilkan gejala yang sama: pengukuran terlihat membaik
        tanpa ada yang membaik.
  \item Untuk fungsi tangga seperti pembulatan, gradien komposisi nol di hampir semua
        titik, sehingga serangan satu langkah menghasilkan $x'=x$ dan ``akurasi di bawah
        serangan'' sama persis dengan akurasi bersih.
  \item BPDA mengganti gradien potongan yang tidak terdiferensiasi dengan gradien
        pengklasifikasi di titik sesudah potongan. Galatnya dibatasi $L\delta$: masuk
        akal ketika penggeseran kecil terhadap skala lekukan, tidak memberi jaminan
        ketika tidak.
  \item EOT memaksimalkan harapan kerugian atas keacakan pertahanan; rata-rata $m$
        sampel tak bias dengan variansi sebanding $1/m$. Keacakan menaikkan biaya per
        langkah, bukan mengubah arah.
  \item Pada model mainan delapan fitur dengan $\varepsilon=0{,}25$, pembulatan
        $q=0{,}01$ menjatuhkan keberhasilan naif dari $17{,}40$ ke $0{,}02$ persen,
        sedangkan BPDA tetap $17{,}40$ persen; $q=0{,}20$ juga tidak menahan BPDA.
  \item Pada NIDS, FGSM menjatuhkan akurasi SVM dari $0{,}998$ ke $0{,}150$, tetapi
        pelatihan ulang harian memulihkannya ke $0{,}998$. Pertahanan yang tahan lama di
        sana datang dari biaya akses dan batasan protokol, bukan dari pengaburan gradien.
\end{itemize}
\end{ringkasan}

\section{Latihan}
\tingkat{menengah}

\begin{latihan}
\begin{enumerate}
  \item Sebuah makalah menulis: ``Pertahanan kami menurunkan keberhasilan serangan dari
        $98$ menjadi $0$ persen pada $\varepsilon=8/255$.'' Sebutkan tiga pemeriksaan
        sebelum menerima kalimat itu, dan apa yang diharapkan pada tiap pemeriksaan bila
        pertahanannya hanya mengaburkan gradien.
  \item Turunkan batas galat BPDA untuk $f(x)=w\cdot x+b$, tunjukkan galatnya nol, lalu
        jelaskan mengapa hasil itu tidak berlaku untuk model dua lapis beraktivasi ReLU.
  \item Pada model mainan bagian empat, tentukan nilai $q$ terkecil yang membuat satu
        langkah BPDA gagal membalik label $x_0$, nyatakan sebagai fungsi margin awal dan
        $\varepsilon$, lalu periksa dengan program.
  \item Anda punya $1\,000$ kueri. Pertahanan mengacak topeng fitur dengan $r=0{,}5$ dan
        menolak menjawab bila lebih dari separuh fitur hilang. Hitung berapa langkah PGD
        yang bisa dijalankan dengan $m=20$ dan dengan $m=1$, lalu jelaskan apa yang Anda
        peroleh dan apa yang Anda bayar.
  \item Ulang percobaan bagian lima dengan model dua lapis beraktivasi ReLU. Laporkan
        akurasi bersih, keberhasilan naif, dan keberhasilan BPDA untuk tiga nilai $q$,
        lalu jelaskan mengapa kosinusnya tidak lagi mendekati satu.
  \item Sebuah NIDS membaca lima puluh fitur, tujuh di antaranya bendera biner dan tiga
        nomor protokol. Hitung berapa fitur yang benar-benar bisa diubah penyerang, sebut
        asumsinya, dan bandingkan dengan gambar $32\times32$ yang punya $3\,072$ nilai.
  \item Jelaskan mengapa pelatihan ulang harian memulihkan akurasi terhadap serangan
        bergradien tetapi tidak terhadap gangguan buta. Kaitkan dengan lemma EOT dan
        definisi serangan adaptif.
  \item Rancang percobaan yang memisahkan dua hipotesis: (a) pertahanan menaikkan biaya
        serangan, (b) pertahanan hanya merusak pengukuran. Sebutkan metrik, jumlah seed,
        bentuk kurva, dan hasil seperti apa yang membuat Anda menerima (a).
\end{enumerate}
\end{latihan}

\peta{Bab 30 pada Athalye dkk., \emph{Obfuscated Gradients Give a False Sense of
Security} (2018): bagian 3 untuk tiga keluarga gradien yang dikaburkan, bagian 4 untuk
BPDA dan serangan adaptif, bagian 6 untuk daftar periksa. Pada Madry dkk., \emph{Towards
Deep Learning Models Resistant to Adversarial Attacks} (2018): bagian 3 untuk PGD dan
titik awal acak. Pada elShehaby dan Matrawy, \emph{Evasion Adversarial Attacks Remain
Impractical Against ML-based Network Intrusion Detection Systems}, untuk pohon serangan,
taksonomi kepraktisan, dan percobaan pelatihan ulang. Pada Bab 20 buku ini untuk aturan
tiga, dan pada Bab 27 untuk FGSM dan PGD.}
